
\exo

Voici les informations fournies par un distributeur de billets concernant les retraits effectués au cours de la semaine dernière :
\begin{center}
\setlength{\tabcolsep}{0.1cm}
\renewcommand{\arraystretch}{1.2}
\begin{tabular}{
>{\raggedright\arraybackslash}m{6cm}
*{7}{>{\centering\arraybackslash}m{1.2cm}}}
\hline 
\textbf{Montant du retrait (en euros)} & $40$ & $50$ & $60$ & $80$ & $100$ & $150$ & $200$ \\ 
\hline
\textbf{Nombre de retraits effectuées} & $55$ & $62$ & $40$ & $23$ & $12$ & $6$ & $2$ \\ 
\hline 
\end{tabular} 
\end{center}
\begin{enumerate}
\item Combien de retraits ont été effectués à ce distributeur la semaine dernière ?
\item Quel est le montant moyen d'un retrait ?
\end{enumerate}

\exo

Pendant huit semaines, Alice a noté chaque matin le temps passé dans les transports en commun pour se rendre à son travail. Les résultats sont données dans le tableau suivant :
\begin{center}
\setlength{\tabcolsep}{0.1cm}
\renewcommand{\arraystretch}{1.2}
\begin{tabular}{
>{\raggedright\arraybackslash}m{4cm}
*{5}{>{\centering\arraybackslash}m{1.2cm}}}
\hline 
\textbf{Temps (en minutes)} & $28$ & $29$ & $33$ & $37$ & $48$ \\ 
\hline
\textbf{Nombre de jours} & $4$ & $8$ & $12$ & $13$ & $3$ \\ 
\hline 
\end{tabular} 
\end{center}
Déterminer la durée moyenne du trajet d'Alice pour se rendre à son travail (à $0,1$ minute près).

\exo

Une entreprise fabrique des éprouvettes en verre, dont le volume théorique est \SI[retain-explicit-plus]{12}{\milli\liter}. 

Sur l'ensemble de la production de l'année dernière, on a prélevé $525$ éprouvettes au hasard, dont les volumes, arrondis à $0,1$ \si{\milli\liter} près, sont donnés dans le tableau suivant :
\begin{center}
\setlength{\tabcolsep}{0.1cm}
\renewcommand{\arraystretch}{1.2}
\begin{tabular}{
>{\raggedright\arraybackslash}m{3.5cm}
*{10}{>{\centering\arraybackslash}m{1.1cm}}}
\hline 
\textbf{Volume (en \si{\milli\liter})} & $11,6$ & $11,7$ & $11,8$ & $11,9$ & $12,0$ & $12,1$ & $12,2$ & $12,3$ & $12,4$ & $12,5$ \\ 
\hline
\textbf{Effectif} & $11$ & $33$ & $60$ & $85$ & $105$ & $100$ & $63$ & $36$ & $23$ & $9$ \\ 
\hline 
\end{tabular} 
\end{center}
Calculer, à \SI[retain-explicit-plus]{0,1}{\milli\liter} près, le volume moyen d'une éprouvette de cet échantillon.

\exo

On considère la série statistique :
\begin{center}
12 -- 9 -- 6 -- 12 -- 10 -- 9 -- 6 -- 9 -- 12 -- 16 -- 9 -- 9.
\end{center}
Regrouper les valeurs de cette série, puis calculer sa moyenne.

\exo

Le professeur de SVT d'Alice fait passer des contrôles qui sont notés sur $100$ points.

Sur ses quatre premiers contrôles, Alice a une moyenne égale à $60$ points, et au cinquième contrôle, elle a obtenu $80$ points. Quelle est sa moyenne, sur les cinq contrôles ?

\exo

Un professeur a noté des devoirs sur $30$ points, et la moyenne de la classe est $m=13,5$.
\begin{enumerate}
\item Sur les copies, le professeur a écrit les notes ramenées sur $20$ points.
\begin{enumerate}
\item Quelle opération a-t-il effectué pour obtenir les notes sur $20$ ?
\item En déduire la moyenne $m'$ des notes sur $20$.
\end{enumerate}
\item Jugeant que cette moyenne est trop faible, le professeur décide d'ajouter $1$ point sur $20$ à tous les élèves. Quelle est alors la nouvelle moyenne $m''$ de la classe ?
\end{enumerate}

\exo

Le tableau suivant donne la répartition des salaires des employés d'une entreprise :
\begin{center}
\setlength{\tabcolsep}{0.1cm}
\renewcommand{\arraystretch}{1.2}
\begin{tabular}{
>{\raggedright\arraybackslash}m{3.5cm}
*{10}{>{\centering\arraybackslash}m{1.1cm}}}
\hline 
\textbf{Salaire (en \euro{})} & \np{1200} & \np{1400} & \np{1600} & \np{2000} & \np{3500} & \np{5000} \\ 
\hline
\textbf{Effectif} & $12$ & $14$ & $13$ & $5$ & $5$ & $1$ \\ 
\hline 
\end{tabular} 
\end{center}
\begin{enumerate}
\item Calculer le salaire moyen dans cette entreprise.
\item Le directeur financier propose d'augmenter tous les salaires de $40$ \euro{}. Quel serait alors le nouveau salaire moyen ?
\item Le P.-D.G. de l'entreprise préfèrerait, lui, augmenter tous les salaires de \SI[retain-explicit-plus]{2}{\percent}. Quel serait alors le nouveau salaire moyen ?
\item Quel est le choix le plus intéressant pour le PDG de l'entreprise ? Et pour les employés dont le salaire est le plus bas ?
\end{enumerate}

\exo

Déterminer la médiane de chacune des séries ci-dessous :
\begin{enumerate}
\item Série \textbf{1} : \quad 4 -- 6 -- 9 -- 11.
\item Série \textbf{2} : \quad 0 -- 5 -- 9 -- 10 -- 17 -- 32.
\item Série \textbf{3} : \quad 10 -- 20 -- 25 -- 30 -- 32 -- 40 -- 55.
\item Série \textbf{4} : \quad 464 -- 631 -- 585 -- 576 -- 539 -- 632 -- 504 -- 601 -- 573 -- 698 -- 493 -- 655 -- 740 -- 671.
\end{enumerate}

\exo

Le responsable de l'association sportive d'un lycée a relevé l'âge des élèves membres de l'association :
\begin{center}
\setlength{\tabcolsep}{0.1cm}
\renewcommand{\arraystretch}{1.2}
\begin{tabular}{
>{\raggedright\arraybackslash}m{3.5cm}
*{10}{>{\centering\arraybackslash}m{1.1cm}}}
\hline 
\textbf{Âge (en années)} & $14$ & $15$ & $16$ & $17$ & $18$ & $19$ & $20$ \\ 
\hline
\textbf{Nombre d'élèves} & $2$ & $21$ & $20$ & $16$ & $18$ & $12$ & $3$ \\ 
\hline 
\end{tabular} 
\end{center}
Déterminer l'âge médian des membres de l'association sportive.

\exo

Les notes de Français des $25$ élèves d'une classe ont été représentées par un diagramme en bâtons :
\begin{center}
\def\xmin{-3}
\def\xmax{21}
\def\ymin{-0.7}
\def\ymax{7}
\def\xunite{0.5cm}
\def\yunite{0.6cm}
\psset{unit=\xunite, xunit=\xunite, yunit=\yunite, algebraic=true, dimen=middle, dotstyle=*, dotsize=3pt 0, linewidth=0.8pt, arrowsize=3pt 2, arrowinset=0.25, comma=true}
\begin{pspicture*}(\xmin,\ymin)(\xmax,\ymax)
%\psgrid[xunit=\xunite, yunit=\yunite, subgriddiv=0, gridlabels=0, gridcolor=lightgray](0,0)(\xmin,\ymin)(\xmax,\ymax)
\psaxes[labelFontSize=\scriptstyle, xAxis=true, yAxis=true, Dx=1, Dy=1, ticksize=-2pt 0, subticks=1]{->}(0,0)(0,0)(\xmax,\ymax)[{\footnotesize Notes},135][{\footnotesize Effectifs},-45]
\multido{\n=1+1}{7}{\psline[linecolor=lightgray, linewidth=0.2pt](0,\n)(\xmax,\n)}
\psline[linewidth=4pt](4,0)(4,1)
\psline[linewidth=4pt](6,0)(6,6)
\psline[linewidth=4pt](7,0)(7,2)
\psline[linewidth=4pt](8,0)(8,5)
\psline[linewidth=4pt](9,0)(9,2)
\psline[linewidth=4pt](10,0)(10,1)
\psline[linewidth=4pt](11,0)(11,1)
\psline[linewidth=4pt](12,0)(12,2)
\psline[linewidth=4pt](13,0)(13,3)
\psline[linewidth=4pt](14,0)(14,1)
\psline[linewidth=4pt](18,0)(18,1)
\end{pspicture*}
\end{center}
Déterminer la médiane de cette série statistique, et en donner une interprétation.

\exo

On a réalisé plusieurs fois l'expérience aléatoire consistant à lancer simultanément trois pièces de monnaie, et à noter le nombre de \og \textsc{pile} \fg{} obtenus. Voici les résultats :
\begin{center}
\setlength{\tabcolsep}{0.1cm}
\renewcommand{\arraystretch}{1.2}
\begin{tabular}{
>{\raggedright\arraybackslash}m{3.5cm}
*{4}{>{\centering\arraybackslash}m{1.1cm}}}
\hline 
\textbf{Nombre de \og \textsc{pile} \fg{} } & $0$ & $1$ & $2$ & $3$ \\ 
\hline
\textbf{Nombre de lancers} & $2$ & $9$ & $11$ & $3$ \\ 
\hline 
\end{tabular} 
\end{center}
Calculer l'effectif total de cette série, puis sa moyenne, et enfin son écart-type (arrondi à $0,1$ près).

\exo

Alice a relevé le DAS (débit d'absorption spécifique, exprimé en \si{\watt\per\kg}) de différents smartphones :
\begin{center}
$0,17$ -- $0,34$ -- $0,79$ -- $0,98$ -- $1,32$.
\end{center}
Calculer l'effectif total de cette série, puis sa moyenne, et enfin son écart-type (arrondi à $0,01$ près).

\exo

Vrai ou faux ? Justifier la réponse.
\begin{enumerate}
\item La moyenne d'une série est toujours une valeur de cette série.
\item Si toutes les valeurs d'une série sont nulles, alors l'écart-type de la série est nul.
\item Si l'écart-type d'une série est nul, alors toutes les valeurs de la série sont nulles.
\item Si on ajoute $1$ à toutes les valeurs d'une série, alors sont écart-type augmente de $1$.
\item Si on divise toutes les valeurs d'une série par $2$, alors son écart-type est divisé par $2$.
\item plus une série comprend de valeurs, plus son écart-type est grand.
\end{enumerate}

\exo

Lors d'une expérience, des élèves ont mesuré l'indice de réfraction d'une plaque de plexiglas :
\begin{center}
$1,47$ -- $1,45$ -- $1,59$ -- $1,46$ -- $1,48$ -- $1,50$ -- $1,46$ -- $1,47$ -- $1,49$ -- $1,47$ -- $1,46$ -- $1,48$.
\end{center}
\begin{enumerate}
\item Déterminer la moyenne et l'écart-type de cette série, arrondis à $0,1$ près.
\item Le professeur estime que la valeur $1,59$ est due à une erreur expérimentale, et décide de la supprimer.
\begin{enumerate}
\item Déterminer la moyenne et l'écart-type de la nouvelle série.
\item Comparer ces résultats avec ceux de la question \textbf{1}.
\end{enumerate}
\end{enumerate}

\exo

Une entreprise est chargée de la maintenance de distributeurs de boissons. Durant une année, le responsable a relevé le nombre d'interventions réalisées sur ces distributeurs. Voici ses résultats :
\begin{center}
\setlength{\tabcolsep}{0.1cm}
\renewcommand{\arraystretch}{1.2}
\begin{tabular}{
>{\raggedright\arraybackslash}m{4.6cm}
*{10}{>{\centering\arraybackslash}m{0.9cm}}}
\hline 
\textbf{Nombre d'interventions} & $1$ & $2$ & $3$ & $4$ & $5$ & $6$ & $7$ & $8$ & $9$ & $10$ \\ 
\hline
\textbf{Nombre de distributeurs} & $10$ & $12$ & $17$ & $44$ & $78$ & $94$ & $83$ & $49$ & $36$ & $16$ \\ 
\hline 
\end{tabular} 
\end{center}
\begin{enumerate}
\item 
\begin{enumerate}
\item Déterminer le nombre moyen $m$ d'interventions réalisées par machine, ainsi que l'écart-type $s$.
\item Le responsable estime qu'il faut changer le modèle de distributeurs si l'intervalle $\intff{m-2s}{m+2s}$ contient moins de \SI[retain-explicit-plus]{95}{\percent} des valeurs de la série. Quelle va être sa décision ?
\end{enumerate}
\item Après vérification, le responsable s'aperçoit qu'il a oublié de compter un distributeur sur lequel on a réalisé trois interventions. Doit-il modifier sa décision ?
\end{enumerate}

\exo

Déterminer les premier et troisième quartiles, ainsi que l'écart interquartile, pour :
\begin{enumerate}
\item les séries définies à l'exercice 8 ;
\item la série définie à l'exercice 9 ;
\item la série définie à l'exercice 10.
\end{enumerate}

\exo

Une municipalité a décidé d'installer des capteurs destinés à relever le niveau sonore dans deux rues de la ville. Ces deux capteurs fournissent chacun $12$ relevés sur une période de $24$ heures. Les mesures sont exprimées en décibels.
\begin{enumerate}
\item Dans la rue de la Liberté, les résultats sont les suivants :
\begin{center}
$55$ -- $50$ -- $52$ -- $56$ -- $64$ -- $74$ -- $79$ -- $65$ -- $73$ -- $74$ -- $64$ -- $50$.
\end{center}
Déterminer la médiane et l'écart interquartile de cette série.
\item Dans la rue de l'Égalité, voici les paramètres de la série obtenue : la médiane vaut $57$, les premier et troisième quartiles valent $55$ et $61$, et les valeurs extrêmes sont $52$ et $64$.

Commenter et comparer la qualité sonore de la vie des habitants de la rue de la Liberté et de la rue de l'Égalité.
\end{enumerate}

